\bookStart{The Norwegian Rune Poem}
\setBookCode{NorwegianRunePoem}

\begin{flushright}%
\textbf{Dating:} Mediæval.%TODO

\textbf{Meter:} Unclear.
\end{flushright}%

\section{Introduction}

The \textbf{Norwegian rune poem} is clearly very closely related to the Icelandic.  With the exception of runes 2 (\emph{úr} ‘slag’) and 4 (\emph{óss} ‘river-mouth’), the names of the runes are identical, as are many of the kennings used to describe them.

Still the language is unmistakably that of mediæval Norway.  As can be seen from the rhymes and alliteration the following uniquely Norwegian sound changes have occurred:
\begin{itemize}
  \item \emph{hl, hn, hr} > \emph{l, n, r} (2 \emph{lęypr} < \emph{hlęypr}; 8 \emph{nęppa} < \emph{hnęppa}; 5 \emph{rossum} < \emph{hrossum}).
  \item \emph{rst} > \emph{st} (5 \emph{vęsta} < \emph{vęrsta})
\end{itemize}

Each rune is described by a couplet of 2 verses linked by end-rhyme and alliteration.  The meter is noticably trochaic, and each verse ends with a rhyming trochée.  The a-verse in 15 out of 16 lines has 2 alliterations; the b-verse always has 1 falling on its very first word.

The description in all cases consists of two clauses or sentences, typically one in each verse.  The first clause gives the name of the rune and a short description of it, while the second clause is more mysterious and often seems to have nothing to do with the name.  A clever solution has however been proposed (TODO: cite), according to which the second clause connects to the visual shape of the rune.  For many of the runes this turns out to be a rewarding line of inquiry, and links between shape and poetry are explored in detail in the notes.

\sectionline

\section{Text}

\bvg\bva[ᚠ]%
\alst{F}é vęldr \alst{f}rę́nda rógi; \hld\ \edtrans{\alst{f}ǿðisk ulfr í skógi}{the wolf rears itself in the wood}{\Bfootnote{Almost certainly visually alluding to the similarity between the rune and the upwards curving branches of the pine tree; the keyword is \emph{skógr} ‘wood, forest’.}}.\eva

\bvb ‘Wealth’ causes the strife of kinsmen; the wolf rears itself in the wood.\evb\evg


\bvg\bva[ᚢ][2]%
\alst{Ú}r er af \alst{ę}ldu \alst{ja}rni; \hld\ \edtrans{\alst{o}pt lęypr ręinn á hjarni}{oft leaps the reindeer on the snow}{\Bfootnote{Perhaps visually alluding to the similarity between the rune and the tracks left by reindeer in the snow; the keyword is \emph{ręinn} ‘reindeer’.}}.\eva

\bvb ‘Slag’ comes from glowing iron; oft leaps the reindeer on the snow.\evb\evg


\bvg\bva[ᚦ][3]%
Þurs vęldr \alst{k}vinna \alst{k}villu; \hld\ \edtrans{\alst{k}átr verðr fár af illu}{few are gladdened by evil}{\Bfootnote{Visually alluding to the similarity between the rune and a smiling face; the keyword is \emph{kátr} ‘glad, happy’.}}.\eva

\bvb ‘Thurse’ causes women’s ailment; few are gladdened by evil.\evb\evg


\bvg\bva[ᚬ][4]%
Óss er \alst{f}lęstra \alst{f}ęrða \hld\ \alst{f}ǫr, en skalpr er sverða.\eva

\bvb ‘River-mouth’ is the start of most journeys, but the scabbard-mouth of swords.\evb\evg


\bvg\bva[ᚱ][5]%
\alst{R}ęið kveða \alst{r}ossum vęsta; \hld\ \alst{R}ęginn sló sverð’it bęsta.\eva

\bvb ‘Chariot’ they say is worst for the horses; Rein forged the best sword.\evb\evg


\bvg\bva[ᚴ][6]%
Kaun er \alst{b}arna \alst{b}ǫlvan; \hld\ \alst{b}ǫl gørvir ná fǫlvan.\eva

\bvb ‘Ulcer’ is the curse of children; a bale makes a corpse pale.\evb\evg


\bvg\bva[ᚼ][7]%
Hagall er \alst{k}aldastr \alst{k}orna; \hld\ \edtrans{\alst{K}ristr skóp hęim’inn forna}{Christ created the world of yore}{\Bfootnote{Visually alluding to the similarity between the rune and the chi-rho; the keyword is \emph{Kristr} ‘Christ’.}}.\eva

\bvb ‘Hail’ is coldest of kernels; Christ created the world of yore.\evb\evg


\bvg\bva[ᚾ][8]%
\alst{N}auðr gørir \alst{n}ęppa kosti; \hld\ \edtrans{\alst{n}øktan kęlr í frosti}{the naked man freezes in the cold}{\Bfootnote{Perhaps visually alluding to the similarity between the rune and the human body in the form of a naked stick-figure; the keyword is \emph{nøktr} ‘naked’.}}.\eva

\bvb ‘Need’ gives scant choice; the naked man freezes in the cold.\evb\evg


\bvg\bva[ᛁ][9]%
Ís kǫllum \alst{b}rú \alst{b}ręiða; \hld\ \edtrans{\alst{b}lindan þarf at lęiða}{the blind man must be lead}{\Bfootnote{Perhaps visually alluding to the similarity between the rune and a blind man’s cane or walking stick; the keyword is \emph{blindr} ‘blind’.}}.\eva

\bvb ‘Ice’ we call a broad bridge; the blind man must be lead.\evb\evg


\bvg\bva[ᛅ][10]%
Ár er \alst{g}umna \alst{g}óði; \hld\ \alst{g}et’k at ǫrr var Fróði.\eva

\bvb ‘Year’ is men’s boon; I recall that Frood was mad.\evb\evg


\bvg\bva[ᛋ][11]%
Sól er \alst{l}anda \alst{l}jómi; \hld\ \edtrans{\alst{l}úti’k hęlgum dómi}{I bow to the holy shrine}{\Bfootnote{Visually alluding to the similarity between the rune and the shape of a kneeling man; the keyword is \emph{lúti’k} ‘I bow’.}}.\eva

\bvb ‘Sun’ is the light of the lands; I bow to the holy shrine.\evb\evg


\bvg\bva[ᛏ][12]%
Týr er \alst{ę}in-ęndr \alst{á}sa; \hld\ \alst{o}pt verðr smiðr blása.\eva

\bvb ‘Tew’ is the one-handed of the Eese; often the smith must blow.\evb\evg


\bvg\bva[ᛒ][13]%
Bjarkan er \alst{l}auf-grǿnstr \alst{l}ima; \hld\ \alst{L}oki bar flę́rða tíma.\eva

\bvb ‘Birch’ is the leaf-greenest of branches; Lock was successful in cunning.\evb\evg


\bvg\bva[ᛘ][14]%
\alst{M}aðr er \alst{m}oldar auki; \hld\ \edtrans{\alst{m}ikil er gręip á hauki}{mighty is the grip on the hawk}{\Bfootnote{Visually alluding to the shape of the bird’s talon; the keyword is \emph{haukr} ‘hawk’.}}.\eva

\bvb ‘Man’ is the product of dust; mighty is the grip on the hawk.\evb\evg


\bvg\bva[ᛚ][15]%
Lǫgr er þar er \alst{f}ęllr ór \alst{f}jalli \hld\ \alst{f}oss; en gull eru nossir.\eva

\bvb ‘Water’ is when from the fell falls a torrent, but gold is treasures.\evb\evg


\bvg\bva[ᛦ][16]%
Ýr er \alst{v}etr-grǿnstr \alst{v}iða; \hld\ \alst{v}ę́nt er, er brennr, at sviða. \eva

\bvb ‘Yew’ is winter-greenest of trees; ’tis expected, when it burns, to get singed.\evb\evg

\sectionline
